\documentclass[10pt,a4paper]{amsart}
\usepackage[margin=19mm]{geometry}
\usepackage{amsmath,amssymb,mathtools}
\usepackage[scr=rsfs,cal=euler]{mathalfa}
\usepackage{xfrac}
\usepackage[unicode,hidelinks,bookmarks=false]{hyperref}
\newcommand{\ie}{i.e.\ }
\newcommand{\cf}{cf.\ }
\newcommand{\resp}{respectively\ }
\newcommand{\Subsection}{\subsection}
\providecommand{\nobreakdash}{\nobreak\hbox{-}}
\setlength{\emergencystretch}{2em}
\multlinegap0pt
\usepackage{tikz}
\usetikzlibrary{cd}
\usepackage{amssymb}
\usepackage{mathtools}
\usepackage[shortlabels]{enumitem}
\usepackage{dsfont}
\usepackage{enumerate}
\usepackage{tikz}
\usepackage{xcolor}
\newtheorem{lemma}{Lemma}
\newtheorem{theorem}{Theorem}
\newcommand{\BF}{\mathcal{BF}}
\newcommand{\Div}{\textup{Div}}
\DeclareMathOperator{\Id}{Id}
\newcommand{\QQ}{\mathbb{Q}}
\newcommand{\Ql}{\QQ_\ell}
\newcommand{\ZZ}{\mathbb{Z}}
\newcommand{\Zl}{\ZZ_\ell}
\DeclareMathOperator{\coker}{coker}
\DeclareMathOperator{\rank}{rank}
\newcommand{\sR}{\mathscr{R}}
\title{Iwasawa theory for abelian towers of digraphs}
\author{Antonio Lei, Katharina M\"uller}
\date{Source: arXiv:2601.19571v1 (CC BY 4.0)}
\begin{document}
\maketitle

\noindent\emph{Source and licence.} Iwasawa theory for abelian towers of digraphs. Version arXiv:2601.19571v1 (CC BY 4.0), distributed by arXiv under the Creative Commons Attribution 4.0 International licence, http://creativecommons.org/licenses/by/4.0/, as recorded in the arXiv licence metadata for this exact version and shown on the abstract page https://arxiv.org/abs/2601.19571v1. Full licence notice: \texttt{licenses/arxiv-2601.19571-CC-BY-4.0.txt}.\par\smallskip

\noindent\textit{Editorial scope.} Counted unit: the article's lemma formula-a,b (source0.tex lines 1028--1035). The article's complete original proof of it is delivered with it (source0.tex lines 1036--1061, one \texttt{proof} environment of 26 lines). No mathematics is written, repaired or re-derived: the statement and the proof are the article's own lines, in the article's own notation, and no retained line is substituted or re-flowed.  Retained uncounted as the source-local context the proof needs: lines 568--572; lines 610--614; lines 764--767.  Nothing was omitted from any retained span, so the package carries no omission record.  No retained line contains a citation, so the wrapper carries no bibliography and no bibliography entry.  The retained lines contain the article's own diagram code, which the wrapper typesets with the standard tikz packages; no image file is delivered and the package declares no asset dependency.  Editorial formatting only, in addition: an AMS \texttt{amsart} preamble that restates the article's own declarations and macros used by the retained lines, namely the environments \texttt{lemma}, \texttt{theorem} on separate counters, together with the standard packages those lines need.  The retained environments print in this document as Lemma 1, Theorem 1 in order of appearance, whereas the article prints the same items with the numbering of its own full document.  The complete native source of the article as distributed in the arXiv e-print for this exact version is bundled unchanged as \texttt{sources/193465/source0.tex}.
\begin{lemma}
\label{isomorphism-group rings}
Let  $\sR_n^{(l)}$ denote the set of classes $[\omega]\in\sR^{(l)}$ such that $\omega\in\Gamma_n^\vee$.     We have the following ring isomorphisms
    \[\Lambda_l(\Gamma)/(T_1^{p^n}-1,T_2^{p^n}-1,\dots , T_d^{p^n}-1)\cong\Zl[\Gamma_n]\cong \bigoplus_{[\omega]\in\sR_n^{(l)}}\Zl[\omega(\Gamma_n)]  \]
   \end{lemma}

\begin{theorem}
\label{Artin-formalism}
   The following equality holds:
    \[Z_{X_n}(u)=\prod_{\omega\in\Gamma_n^\vee}L(X_n/X,\omega,u).\]
\end{theorem}

\begin{theorem}
\label{char-udir}
Let $r$ be the order of $X_0$. The $\Lambda_l(\Gamma)$-modules ${\rm Pic}_{l}(X_{\infty})$ and $\BF_l(X_\infty)$ are isomorphic to the quotients $\Lambda^r/\left(D_X-A_\alpha\right)$ and $\Lambda^r/\left(\Id-A_\alpha\right)$, respectively.
\end{theorem}

\begin{lemma}
\label{formula-a,b}
    We have the following identities
    \begin{align*}
        a(X_n)&=\sum_{\omega\in \Gamma_n^\vee}a(X_n,\omega),\\
        b(X_n)&=\sum_{\omega\in \Gamma_n^\vee}b(X_n,\omega).
    \end{align*}
\end{lemma}

\begin{proof}
    The equation for the analytic ranks follows from the Artin formalism (Theorem \ref{Artin-formalism}). It remains to prove the equation for the algebraic ranks. 

   Let $l$ be a prime number distinct from $p$. For each $\omega\in\Gamma_n^\vee$, let $f_\omega$ denote $[\Ql(\omega(\Gamma_n)):\Ql]$. There are $f_\omega$ characters in the equivalence class $[\omega]$ and
   \[
   b(X_n,\omega)=b(X_n,\omega')\ \forall \omega'\in[\omega],
   \]
   and  
   \begin{equation}
   \label{eq:b-formula}\sum_{\omega'\in[\omega]}b(X_n,\omega')=f_\omega\cdot  b(X_n,\omega)=\rank_{\Zl}\Zl[\omega(\Gamma_n)]^r/(\mathrm{Id}-\omega(A_\alpha)).
   \end{equation}


 Recall from Lemma~\ref{isomorphism-group rings} that
    \[ \varphi_n:\Zl[\Gamma_n]\cong\bigoplus_{[\omega]\in\sR_n^{(l)}}\Zl[\omega(\Gamma_n)].\]
Further, we have seen in the proof of Theorem~\ref{char-udir} that there is an isomorphism $\Div_l(X_n)\cong\displaystyle \bigoplus_{v\in V(X)}\Zl[\Gamma_n](v,1)$, which gives
    \[\BF_l(X_n)\cong \Zl[\Gamma_n]^r/(\textup{Id}-A_\alpha).\] 
    
    For each $[\omega]\in\sR_n^{(l)}$,  define the map $\psi_{[\omega]}$ on $\Zl[\omega(\Gamma_n)]^r$ given by the matrix $\textup{Id}-\omega(A_\alpha)$. Let $\psi_n=\displaystyle\bigoplus_{[\omega]\in\sR_n^{(l)}}\psi_{[\omega]}$ be the corresponding map on $\displaystyle\bigoplus_{m\le n}\Zl[\omega_m(\Gamma_n)]^r$. We have the following commutative diagram
    \[\begin{tikzcd}
        \Zl[\Gamma_n]^r\arrow[r,"\textup{Id}-A_\alpha"]\arrow[d,"\varphi_n"]&\Zl[\Gamma_n]^r\arrow[d,"\varphi_n"]\\
        \left(\bigoplus_{m\le n}\Zl[\omega_m(\Gamma_n)]\right)^r\arrow[r, "\psi_n"]&\left(\bigoplus_{m\le n}\Zl[\omega_m(\Gamma_n)]\right)^r.
    \end{tikzcd}\]
    In particular, \[b(X_n)=\rank_{\Zl}(\coker(\psi_n))=\sum_{[\omega]\in \sR_n^{(l)}}\rank_{\Zl}(\coker(\psi_{[\omega]}))=\sum_{\omega\in \Gamma_n^\vee}b(X_n,\omega),\]
    where the last equality follows from \eqref{eq:b-formula}.
\end{proof}

\end{document}
